\documentclass[11pt,fontset=fandol]{ctexart}
\usepackage[a4paper,margin=23mm]{geometry}
\usepackage{amsmath,amssymb,mathtools,bm,mathrsfs,esint}
\usepackage{graphicx,longtable,booktabs,array,calc}
\usepackage[normalem]{ulem}
\usepackage{xcolor,hyperref}
\hypersetup{hidelinks}
\setcounter{secnumdepth}{-1}
\newcounter{none}
\providecommand{\tightlist}{\setlength{\itemsep}{0pt}\setlength{\parskip}{0pt}}
\providecommand{\passthrough}[1]{#1}
\usepackage{amssymb,mathrsfs,esint,bm,newunicodechar}
\allowdisplaybreaks
\setlength{\emergencystretch}{3em}
\providecommand{\llangle}{\langle\!\langle}
\providecommand{\rrangle}{\rangle\!\rangle}
\newunicodechar{∞}{\ensuremath{\infty}}
\newunicodechar{∈}{\ensuremath{\in}}
\newunicodechar{∑}{\ensuremath{\sum}}
\newunicodechar{∫}{\ensuremath{\int}}
\newunicodechar{≤}{\ensuremath{\leq}}
\newunicodechar{≥}{\ensuremath{\geq}}
\newunicodechar{≠}{\ensuremath{\neq}}
\newunicodechar{→}{\ensuremath{\to}}
\newunicodechar{⇒}{\ensuremath{\Rightarrow}}
\newunicodechar{α}{\ensuremath{\alpha}}
\newunicodechar{β}{\ensuremath{\beta}}
\newunicodechar{γ}{\ensuremath{\gamma}}
\newunicodechar{λ}{\ensuremath{\lambda}}
\newunicodechar{φ}{\ensuremath{\varphi}}
\newunicodechar{ε}{\ensuremath{\varepsilon}}
\newunicodechar{σ}{\ensuremath{\sigma}}
\newunicodechar{θ}{\ensuremath{\theta}}
\newunicodechar{δ}{\ensuremath{\delta}}
\newunicodechar{η}{\ensuremath{\eta}}
\newunicodechar{π}{\ensuremath{\pi}}
\newunicodechar{ξ}{\ensuremath{\xi}}
\newunicodechar{Σ}{\ensuremath{\Sigma}}
\newunicodechar{Δ}{\ensuremath{\Delta}}
\newunicodechar{∂}{\ensuremath{\partial}}
\newunicodechar{∀}{\ensuremath{\forall}}
\newunicodechar{∃}{\ensuremath{\exists}}
\newunicodechar{√}{\ensuremath{\surd}}
\newunicodechar{−}{\ensuremath{-}}
\newunicodechar{μ}{\ensuremath{\mu}}
\newunicodechar{ρ}{\ensuremath{\rho}}
\newunicodechar{τ}{\ensuremath{\tau}}
\newunicodechar{ω}{\ensuremath{\omega}}
\newunicodechar{Ω}{\ensuremath{\Omega}}
\newunicodechar{χ}{\ensuremath{\chi}}
\newunicodechar{ψ}{\ensuremath{\psi}}
\newunicodechar{Γ}{\ensuremath{\Gamma}}
\newunicodechar{Φ}{\ensuremath{\Phi}}
\newunicodechar{Λ}{\ensuremath{\Lambda}}
\newunicodechar{ℝ}{\ensuremath{\mathbb{R}}}
\newunicodechar{ℤ}{\ensuremath{\mathbb{Z}}}
\newunicodechar{ℕ}{\ensuremath{\mathbb{N}}}
\newunicodechar{△}{\ensuremath{\triangle}}
\newunicodechar{⊗}{\ensuremath{\otimes}}
\begin{document}
\section{2017年杭州电子科技大学601数学分析考研真题}\label{ux5e74ux676dux5ddeux7535ux5b50ux79d1ux6280ux5927ux5b66601ux6570ux5b66ux5206ux6790ux8003ux7814ux771fux9898}

来源：\texttt{HDU/601数学分析/2017/2017年杭州电子科技大学601数学分析考研真题.pdf}。原
PDF 页序：1, 2。

\section{杭州电子科技大学2017年攻读硕士学位研究生招生考试}\label{ux676dux5ddeux7535ux5b50ux79d1ux6280ux5927ux5b662017ux5e74ux653bux8bfbux7855ux58ebux5b66ux4f4dux7814ux7a76ux751fux62dbux751fux8003ux8bd5}

\section{《数学分析》试题}\label{ux6570ux5b66ux5206ux6790ux8bd5ux9898}

（试题共四大题，共2页，总分150分）

姓名 \_\_\_\_　报考专业 \_\_\_\_　准考证号 \_\_\_\_

【所有答案必须写在答题纸上，做在试卷或草稿纸上无效！】

\begin{quote}
助手整理依据：交叉读取远程OCR与本地逐页OCR，恢复曲面积分、闭曲线符号和数学排版。一5按两份打印题面保留为求
\(\partial z/\partial y\)，不沿用原参考答案中的
\(\partial z/\partial x\) 作为题目所求。
\end{quote}

\subsection{一、填空题（8小题，每题5分，共40分）}\label{ux4e00ux586bux7a7aux98988ux5c0fux9898ux6bcfux98985ux5206ux517140ux5206}

\begin{enumerate}
\def\labelenumi{\arabic{enumi}.}
\item
  \(\displaystyle\lim_{x\to0}\frac{\ln(1+x^2)}{x\tan2x}=\) \_\_\_\_。
\item
  函数 \[
  f(x)=\begin{cases}e^{1/(x-1)},&x>0,\\x\sin(1/x),&-1<x\le0\end{cases}
  \] 的跳跃间断点为 \(x=\) \_\_\_\_。

  原题第二支在 \(x=0\) 的表达式含 \(1/x\)，第一支在 \(x=1\)
  也无定义；题面按原印保留，其单侧极限判别见解析。
\item
  \(\displaystyle\int_{-1}^1[\ln(1+x^4)\sin^7x+|x|e^{|x|}]\,dx=\)
  \_\_\_\_。
\item
  设 \(\displaystyle\int_0^tf(u)\,du=f(t)-1\)，则 \(f^{(n)}(t)=\)
  \_\_\_\_。
\item
  设 \(e^{-xy}+2z-e^z=0\)，则 \(\partial z/\partial y=\) \_\_\_\_。
\item
  设 \(f(x,y)\) 连续，交换
  \(\displaystyle\int_0^1dx\int_{x^2}^{2-x}f(x,y)\,dy\) 的积分次序，得到
  \_\_\_\_。
\item
  幂级数 \(\displaystyle\sum_{n=1}^\infty\frac{9^n+(-4)^n}{n}x^{2n}\)
  的收敛半径 \(R=\) \_\_\_\_。
\item
  已知 \(z=\arctan(x/y)\)，则 \(dz=\) \_\_\_\_。
\end{enumerate}

\subsection{二、计算题（4小题，共30分）}\label{ux4e8cux8ba1ux7b97ux98984ux5c0fux9898ux517130ux5206}

\begin{enumerate}
\def\labelenumi{\arabic{enumi}.}
\item
  （5分）求 \(\displaystyle\lim_{x\to0}\frac{2^x+3^x-2}{\arcsin x}\)。
\item
  （5分）求 \(\displaystyle\int x^2\arctan x\,dx\)。
\item
  （10分）计算反常积分
  \(\displaystyle\int_0^1\frac{x^5}{\sqrt{1-x^2}}\,dx\)。
\item
  （10分）计算 \[
  \iint_\Sigma x^3\,dy\,dz+y^3\,dz\,dx+z^3\,dx\,dy,
  \] 其中 \(\Sigma\) 是上半球面 \(z=\sqrt{a^2-x^2-y^2}\) 的上侧。
\end{enumerate}

\subsection{三、综合题（4小题，每题10分，共40分）}\label{ux4e09ux7efcux5408ux98984ux5c0fux9898ux6bcfux989810ux5206ux517140ux5206}

\begin{enumerate}
\def\labelenumi{\arabic{enumi}.}
\item
  判别
  \(\displaystyle\sum_{n=1}^\infty\frac{(-1)^nx\arctan x}{1+n^3x^2}\) 在
  \(\mathbb R\) 上的一致收敛性。
\item
  设 \(p>0\)，讨论幂级数
  \(\displaystyle\sum_{n=1}^\infty\frac{x^n}{(n+p)^p}\) 的收敛域。
\item
  已知 \(f\) 在 \(\mathbb R\) 上可导，\(\lim_{x\to\infty}f'(x)=e\)，且
  \[
  \lim_{x\to\infty}\left(\frac{x+k}{x-k}\right)^x
  =\lim_{x\to\infty}[f(x+1)-f(x-1)],
  \] 求常数 \(k\)。
\item
  设 \(a>0\)，求 \(f(x,y,z)=xy^2z^3\)（\(x,y,z>0\)）在 \(x+y+z=a\)
  下的极值，并证明 \[
  xy^2z^3\le108\left(\frac{x+y+z}{6}\right)^6.
  \]
\end{enumerate}

\subsection{四、证明题（4小题，每题10分，共40分）}\label{ux56dbux8bc1ux660eux98984ux5c0fux9898ux6bcfux989810ux5206ux517140ux5206}

\begin{enumerate}
\def\labelenumi{\arabic{enumi}.}
\item
  当 \(0<x<\pi\) 时，证明 \(\sin(x/2)>x/\pi\)。
\item
  设 \(f''(x)<0,f(0)=0\)，证明任意 \(x_1,x_2>0\) 都有
  \(f(x_1+x_2)<f(x_1)+f(x_2)\)。
\item
  证明 \(f(x)=x^2\) 在 \((0,\infty)\) 连续但不一致连续。
\item
  证明： \[
  \begin{aligned}
  \text{（1）}\quad&\oint_Lxe^{\sin y}\,dy-ye^{-\sin x}\,dx
  =\oint_Lxe^{-\sin y}\,dy-ye^{\sin x}\,dx,\\
  \text{（2）}\quad&\oint_Lxe^{\sin y}\,dy-ye^{-\sin x}\,dx\ge2\pi^2,
  \end{aligned}
  \] 其中 \(L\) 为正方形 \(D=\{(x,y):0\le x\le\pi,0\le y\le\pi\}\)
  的逆时针边界。
\end{enumerate}
\end{document}
